\subsection{An evaluated arithmetic parameter and its smaller obstruction map}

\begin{statement}{D.56}{Proved}{The actual order-\(3888\) quotient, its unique arithmetic parameter, and the forty-dimensional correction space}
Retain the actual \(h,f_h,G_3,A,R_h,\alpha_0\)
of D.46--D.54. All source kernels below are closed.
Put \(A_*=\mathbb F_3^4\), with column-vector convention, and set
\[
 \begin{gathered}
 S_*=
 \begin{pmatrix}
 0&1&0&0\\1&0&0&0\\0&0&0&1\\0&0&1&0
 \end{pmatrix},\qquad
 R_*=
 \begin{pmatrix}
 1&0&0&0\\2&1&2&0\\0&0&1&0\\2&0&1&1
 \end{pmatrix},\\
 v_*=(1,1,1,0)^{\mathsf t},\qquad
 Q_*=\langle S_*,R_*\rangle,\qquad
 G_*=A_*\rtimes Q_* .
 \end{gathered}
 \tag{D.56.1}
\]
The following is a specified epimorphism from the original marked source:
\[
 \begin{aligned}
 \varphi_*:\widehat B&\longrightarrow G_*,
 &s&\longmapsto(0,S_*),
 &r&\longmapsto(v_*,R_*).
 \end{aligned}
 \tag{D.56.2}
\]
It factors through the actual \(G_3=B/(N'N^3)\).
The induced linear quotient and its kernel satisfy
\[
 \begin{gathered}
 Q\twoheadrightarrow Q_*,
 \qquad |Q_*|=48,\qquad
 \ker(Q\to Q_*)\simeq A_4,\qquad
 5I\longmapsto-I_4,\\
 |G_*|=3888,\qquad
 1\longrightarrow C_3^{\,93}\longrightarrow
 P_*\longrightarrow A_4\longrightarrow1,
 \quad P_*=\ker(G_3\to G_*),\\
 |P_*|=12\cdot3^{93},\qquad H_1(P_*,\mathbb F_2)=0 .
 \end{gathered}
 \tag{D.56.3}
\]
The \(Q_*\)-module \(A_*\) is irreducible, with
\(\operatorname{End}_{Q_*}(A_*)=\mathbb F_9\).
The image of the actual arithmetic parameter is determined exactly:
\[
 \begin{gathered}
 a_*=\varphi_*(f_h)=(2,1,2,1)^{\mathsf t},\qquad
 \varphi_*(A)=A_*,\qquad \varphi_*(C_h)=A_*,\\
 t_*=\varphi_*(sr),\qquad
 |t_*|=12,\qquad
 t_*^6=((0,2,1,0)^{\mathsf t},-I_4),\\
 \varphi_*\alpha_h=\operatorname{Int}(t_*^6)\varphi_* .
 \end{gathered}
 \tag{D.56.4}
\]
Here \(\operatorname{Int}(k)(g)=kgk^{-1}\).
In particular the normalized arithmetic action preserves this
specific quotient and acts as \(-I_4\) on its translation module.
The map from the spliced D.49 source factors through its identity
three-primary component:
\[
 \mathcal W_h\longrightarrow P_{h,3}
       \longrightarrow G_*,
 \qquad \alpha_0|_{G_*}=\operatorname{id}.
 \tag{D.56.5}
\]
These are the two actual comparison maps for the original action
and its retained splice.

Write \(K_*=\ker(\widehat B\to G_*)\), and put
\[
 \begin{gathered}
 U_*=K_*^{\mathrm{ab},(2)},\quad V_*=U_*/2,\quad
 e_*={1\over81}\sum_{a\in A_*}a,\quad
 U_*^-=(1-e_*)U_*,\quad W_*=U_*^-/2 .
 \end{gathered}
 \tag{D.56.6}
\]
The actual marked graph and its compactification give
\[
 \begin{array}{c|r}
 \text{quantity}&\text{value}\\ \hline
 \operatorname{rank}_{\mathbb Z_2}U_*&649\\
 \operatorname{rank}_{\mathbb Z_2}e_*U_*&9\\
 \dim W_*&640\\
 \text{cusp count}&324\\
 \text{cusp radical dimension}&323\\
 \text{compact homology dimension}&326\\
 \text{moving cusp radical dimension}&320\\
 \text{moving compact homology dimension}&320
 \end{array}.
 \tag{D.56.7}
\]
Every dimension in the last six rows is over \(\mathbb F_2\).
The genus is \(163\). The peripheral maps are the original
right multiplications by \(\varphi_*(s),\varphi_*(r)\);
the cusp loops are the conjugates of \(t^{12}\).

The actual correction and same-action parameter maps are
\[
 \begin{gathered}
 \mathcal P_*=\ker\!\left[
 W_*^2\xrightarrow{\,L_*\,}W_*\right],
 \qquad L_*(b,c)=(1-s)b+s(1-r)c,\\
 \mathcal K_*=
 \{(u_s+d,u_r+d):
       u_s\in W_*^s,\ u_r\in W_*^r,\ d\in W_*^t\},
 \qquad \mathscr D_*=\mathcal P_*/\mathcal K_*,\\
 \begin{array}{c|rrrrrr}
 &W_*^s&W_*^r&W_*^t&
       \mathcal P_*&\mathcal K_*&\mathscr D_*\\ \hline
 \dim&336&210&54&640&600&40
 \end{array}.
 \end{gathered}
 \tag{D.56.8}
\]
The map defining \(\mathcal K_*\) is injective.
Thus \(640\) counts actual peripheral conjugator-pair parameters,
and \(40\) counts their induced homology derivations.

There are surjective maps from the full arithmetic modules and
parameters, all induced by \(K_3\hookrightarrow K_*\):
\[
 \begin{gathered}
 u:U\twoheadrightarrow U_*,
 \quad j:V\twoheadrightarrow V_*,
 \quad j_W:W\twoheadrightarrow W_*,
 \quad \xi:\mathcal P\twoheadrightarrow\mathcal P_*,
 \quad \mathscr D\twoheadrightarrow\mathscr D_*,\\
 \theta:Z=W\wedge V\twoheadrightarrow
       Z_*=W_*\wedge V_*,
 \qquad \tau:Z_{G_3}\twoheadrightarrow (Z_*)_{G_*},\\
 \dim Z_*=210240,\qquad
 \dim Z_*^{A_*}=2560,\qquad
 \dim(1-e_*)Z_*=207680 .
 \end{gathered}
 \tag{D.56.9}
\]
For example \(\xi(b,c)=(j_Wb,j_Wc)\), and
\(\theta(w\wedge v)=j_Ww\wedge jv\).
The original chosen D.50 action and its conjugators induce
a specific actual action \(\alpha_{0,*}\) on the smaller
joint homology source. With these choices the actual boundary maps
commute:
\[
 \begin{gathered}
 \tau\bigl(q_{\alpha_0}(p)\bigr)
       =q_{\alpha_{0,*}}\bigl(\xi(p)\bigr),\\
 \overline\tau:
 Z_{G_3}/\ell(\mathcal K)
       \twoheadrightarrow
 (Z_*)_{G_*}/\ell_*(\mathcal K_*),\\
 \overline\tau\,\overline q_{\alpha_0}
       =\overline q_{\alpha_{0,*}}\,
          (\mathscr D\longrightarrow\mathscr D_*).
 \end{gathered}
 \tag{D.56.10}
\]
The right-hand quadratic map is an actual forty-parameter
arithmetic obstruction. Its coefficients and zero locus have
not been evaluated here.

For completeness, its source and projection kernels are explicit.
Set
\[
 \begin{gathered}
 M_*=\ker(K_*\to U_*),\quad
 R_{50,*}=R_h\cap M_*,\quad
 L_{2,*}=\overline{\Gamma_3K_*(K_*')^2K_*^4},\\
 \mathcal Z_{h,*}=\widehat B/R_{50,*},\qquad
 \mathcal T_{2,*}=\widehat B/L_{2,*},\qquad
 \mathcal X_*=\widehat B/(R_{50,*}\cap L_{2,*}).
 \end{gathered}
 \tag{D.56.11}
\]
The image of \(R_h\) in \(\widehat B/M_*\) is exactly \(U_*^-\).
The image of \(R_{50,*}\) in \(\mathcal T_{2,*}\) is exactly \(Z_*\).
Consequently
\[
 \begin{gathered}
 \mathcal X_*=
 \mathcal Z_{h,*}\times_{\mathcal T_{2,*}/Z_*}\mathcal T_{2,*},
 \qquad
 \ker(\mathcal Z_{h,*}\to\mathcal W_h)=U_*^-,\\
 \ker(\mathcal X_*\to\mathcal Z_{h,*})=Z_*,
 \qquad
 \ker(\mathcal X_*\to\mathcal T_{2,*})
       =L_{2,*}R_{50,*}/R_{50,*},\\
 \ker(\mathcal Z_h\to\mathcal Z_{h,*})
       ={R_h\cap M_*\over R_h\cap M_\infty},
 \qquad
 \ker(\mathcal X_{54}\to\mathcal X_*)
       ={R_{50,*}\cap L_{2,*}\over R_{50}\cap L_2}.
 \end{gathered}
 \tag{D.56.12}
\]
Here \(M_\infty=\ker(K_3\to U)\) is D.50's
\(\mathfrak K_\infty\). All maps retain the D.49
parameter and its original conjugator framing.
The ordinary group norms have the further exact comparison
\[
 \theta\!\left(\sum_{g\in G_3}gz\right)
   =|P_*|\sum_{g_*\in G_*}g_*\theta(z)=0
       \quad(z\in Z).
 \tag{D.56.13}
\]
The smaller boundary equation in (D.56.10) therefore remains
part of the comparison, even after the compact norm of D.55
has been calculated.

The certificates use the actual \(Q\)-graph and the actual
arithmetic \(f_h\). They do not substitute the separate
order-\(1080\) family for this quotient.
The manuscript passages are Part D(e), section 51.16,
folio 734 verso, scan 174, and section 51.15,
folios 725--727, scans 156--160.
\end{statement}
\begin{proof}
\emph{The literal finite source.}
Multiplication in \(A_*\rtimes Q_*\) is
\((v,M)(w,N)=(v+Mw,MN)\).
The displayed matrices have orders \(2,3\), and
\((1+R_*+R_*^2)v_*=0\).
Thus (D.56.2) is already a homomorphism from
\(C_2*C_3\), and it extends continuously to its
profinite completion.

There is a particularly short complete certificate that the
linear map factors through the marked \(Q\).
Starting at its identity, assign to a \(Q\)-state the product
of \(S_*,R_*\) along its shortest marked word.
The program checks consistency on every right-\(s\) and
right-\(r\) edge of all \(576\) states.
Every word with the same \(Q\)-endpoint therefore has the
same linear matrix. This proves the claimed quotient map,
rather than only a match of generator orders.
Its complete image has \(48\) matrices. Its kernel has
twelve states: the four marked states
\[
            0,\ 572,\ 574,\ 575
 \tag{D.56.14}
\]
form a Klein four group, and the marked state \(321\)
has order three and cyclically permutes its three nonidentity
elements. The twelve products are all distinct.
This gives the explicit semidirect-product identification
of that kernel with \(A_4\).
The numbering is the original \(Q\)-matrix ordering in
the D.53 certificate; both the matrices and the state maps
are retained, so the identification has a literal inverse.

Since \(N\) has trivial linear image, its affine image is
contained in \(A_*\). Reconstruct the \(97\) marked free
generators of \(N\) from the same spanning tree as D.53.
Their affine images are the columns of the retained
rank-four map
\[
         \phi:N^{\rm ab}/3\twoheadrightarrow A_* .
 \tag{D.56.15}
\]
Every one of the \(97\) word images is checked against
the corresponding collapsed graph cycle. This proves
surjectivity and the exact coordinate map.
In particular \(N'N^3\) is killed. The full affine map
retaining \(Q\) has image \(A_*\rtimes Q\), of order
\(46656\); projection of its linear factor gives (D.56.2).
The \(3888\) actual source states are all generated by
the two displayed images.

The kernel of \(\phi\) has dimension \(93\).
For any of the twelve linear kernel states, lift it to
\(G_3\) and remove its translation using \(\phi(N)\).
This gives the surjection \(P_*\to A_4\), whose kernel
is precisely \(\ker\phi\). It proves the exact sequence
in (D.56.3), including its maps.
The odd kernel has no positive-degree homology with
\(\mathbb F_2\) coefficients, so
\(H_1(P_*,\mathbb F_2)=H_1(A_4,\mathbb F_2)=0\).
Equivalently, \(A_4\) is generated by its elements of
order three, which all map trivially to an abelian
two-group.

The forty projective lines of \(A_*\) are tested directly:
the \(S_*,R_*\)-span of each has dimension four.
This proves irreducibility. Solving the sixteen-entry
matrix commutation equations gives a two-dimensional
\(\mathbb F_3\)-algebra. All eight nonzero elements
have inverses in the same algebra, proving that it is
\(\mathbb F_9\).
The dual character pairs have orbit sizes
\(8,12,12,8\), with stabilizer orders \(6,4,4,6\);
their complete maps are also retained for later
coefficient calculations.

\emph{Determining the actual arithmetic parameter.}
D.49 supplies the original, normalized identities
\[
       \alpha_h(s)=s f_h,\qquad
       \alpha_h(r)=f_h^{-1}r,\qquad
       \alpha_h(t)=t,\qquad f_h\in N,
 \tag{D.56.16}
\]
and D.47 identifies the original \(r\)-conjugator's
\(Q\)-coordinate as \(5I\).
In our module this coordinate is \(-I_4\).
Write \(a=\varphi_*(f_h)\).
The order-two relation for \(s f_h\) gives
\[
          S_*a=-a .
 \tag{D.56.17}
\]
Every affine lift of the original \(r\)-conjugator
has form \((b,-I_4)\), and its square is the identity.
Its conjugation of \((v_*,R_*)\) has translation
\((1-R_*)b-v_*\).
Comparison with \(f_h^{-1}r\) consequently gives
\[
              a-2v_*\in\operatorname{im}(1-R_*) .
 \tag{D.56.18}
\]
These two conditions determine \(a\) without a choice
of its remaining conjugator coordinate.
Indeed \(\operatorname{im}(1-R_*)\) is exactly the
span of the second and fourth coordinate vectors.
Thus the first and third coordinates of \(a\) equal
\(2\), and (D.56.17) makes the other two equal \(1\).
This proves the value in (D.56.4).
It also proves directly why no lift of \(5I\)
centralizes the actual \(r\)-image: that would require
\(2v_*\in\operatorname{im}(1-R_*)\).

The normal span of this nonzero vector is all of the
irreducible \(A_*\). This is the image of both the
actual \(A\) and the actual closed normal subgroup
\(C_h\). Direct multiplication gives the displayed
value of \(t_*^6\), and conjugation by that value sends
the two marked generators to
\[
               s_*a_*,\qquad a_*^{-1}r_* .
 \tag{D.56.19}
\]
Equation (D.56.16) proves the equality of homomorphisms
in (D.56.4). It proves preservation of the quotient
kernel as well as the resulting automorphism.
The certificate independently checks the map on all
\(7776\) directed marked edges and all \(3888\) states.

Finally \(C_h\) has three-group image \(A_*\) in \(G_*\).
The restriction of this finite map to \(C_h\) factors
through \(C_h^{(3)}\). Therefore
\(R_{h,3}\) lies in its source kernel, which proves
the actual factorization through \(P_{h,3}\).
The retained splice is the identity on that factor,
proving (D.56.5). No change of the marked source
or of the original arithmetic parameter is used.

\emph{The smaller homology and its peripheral parameter space.}
The original \(s,r\) orders remain \(2,3\), so \(K_*\)
is torsion-free. Its Bass--Serre graph has \(3888\)
edges and \(1944+1296\) vertices. Its rank is \(649\).
The \(t_*\)-cycles have length \(12\), hence there are
\(324\) cusps and genus
\((649-324+1)/2=163\).
The integral cusp sublattice is primitive, with its
single sum relation. Its modulo-two rank is \(323\).

The coefficient-two averaging idempotent \(e_*\)
is integral because \(81\) is odd.
The quotient graph for \(Q_*\) has rank
\(1+48/6=9\). Its cusp count is \(48/12=4\).
The covering to that graph has group \(A_*\);
both cusp orders are \(12\), so no additional
cusp ramification occurs.
Averaging identifies invariant homology with the
rank-nine quotient homology. It identifies the
invariant cusp radical with its rank-three radical.
The complementary ranks are therefore \(640\)
and \(320\), and the moving compact rank is \(320\).
This proves all the maps and numbers in (D.56.7).

The certificate reconstructs all \(649\) free words
and their integral cycles. It computes the integral
ribbon pairing as one half of the sum of the oriented
trivalent-vertex determinants. Every entry is integral;
its modulo-two rank is \(326\), and its radical is
exactly the span of the \(324\) marked cusp loops.
The source permutations give the homology matrices.
All deck equivariances of this pairing are checked.
The \(A_*\)-norm of each cycle gives the actual
invariant projection; subtracting it gives \(W_*\).
The fixed-space nullities of its \(s,r,t\) matrices
are \(336,210,54\), respectively.
Every matrix and each fixed-space basis is included
in the certificate.

Here \((W_*)_{G_*}=0\), because already
\((W_*)_{A_*}=0\). The image of \(L_*\) is the
augmentation subspace, so \(L_*\) is onto.
Its kernel has dimension \(640\).
The cohomological and peripheral argument of D.53
applies with the literal smaller coefficient map:
the averaging idempotent kills all cohomology of
the \(G_*\)-module \(W_*\), and the two marked
cyclic restrictions identify the kernel of the
pair-to-derivation map with
\[
             W_*^s\oplus W_*^r\oplus W_*^t
        \longrightarrow W_*^2,\quad
             (u_s,u_r,d)\longmapsto(u_s+d,u_r+d).
 \tag{D.56.20}
\]
For its injectivity, a vector in the kernel must be
fixed by \(s,r,t\), and hence by \(G_*\); its
\(A_*\)-invariant part in \(W_*\) is zero.
The dimensions in (D.56.8) now follow.

\emph{Surjective maps from the full arithmetic obstruction.}
The inclusion \(K_3\hookrightarrow K_*\) has quotient
\(P_*\). The image of \(U\to U_*\) has quotient equal
to the maximal abelian pro-two quotient of \(P_*\).
That quotient is trivial by (D.56.3); thus \(u\)
and its reduction \(j\) are onto.
Conjugation by \(\widehat B\) commutes with these maps.
The action on the smaller target factors through
\(G_*\).
Since \(A\to A_*\) is onto and has odd kernel, their
averaging idempotents satisfy
\[
                       u e_A=e_*u .
 \tag{D.56.21}
\]
It follows that \(U^-\to U_*^-\) and \(j_W\) are onto.
Taking exterior products gives \(\theta\);
taking coinvariants gives \(\tau\), with
\[
       \ker\tau
           =\operatorname{im}\bigl((\ker\theta)_{G_3}
                                    \longrightarrow Z_{G_3}\bigr).
 \tag{D.56.22}
\]
This is right exactness of the specified
coinvariant sequence, with the action on \(Z_*\)
inflated along \(G_3\to G_*\).

For surjectivity on peripheral pairs, put
\(J=\ker j_W\). Exactness of averaging implies
\(J_A=0\). Therefore
\[
     (1-s)J+s(1-r)J=J .
 \tag{D.56.23}
\]
Lift any pair in \(\mathcal P_*\) to \(W^2\).
Its cusp defect belongs to \(J\); (D.56.23)
removes it using a pair in \(J^2\).
This proves surjectivity of \(\xi\).
The image of \(\mathcal K\) lies in \(\mathcal K_*\),
so it also proves surjectivity
\(\mathscr D\to\mathscr D_*\), with literal kernel
\(\xi^{-1}(\mathcal K_*)/\mathcal K\).

As an \(A_*\)-module, \(W_*\) consists of eight
copies of its augmentation module.
Indeed the graph chain modules are free
\(\mathbb F_2[A_*]\)-modules, since \(A_*\) has
trivial intersection with the two vertex stabilizers.
Their moving chain difference has multiplicity
eight; semisimplicity over the odd group \(A_*\)
gives the stated isomorphism.
The eighty nontrivial characters occur in forty
inverse pairs, each with multiplicity eight.
Thus
\(\dim(\bigwedge^2W_*)^{A_*}=40\cdot8^2=2560\).
The cross term with \(V_*^{A_*}\) has no invariant
part. Also
\[
 \dim(W_*\wedge V_*)
    ={640\choose2}+640\cdot9=210240 .
 \tag{D.56.24}
\]
Subtracting the invariant dimension gives \(207680\).
Exactness of odd averaging also proves that
\(Z^A\to Z_*^{A_*}\) is onto.

\emph{The actual joint source and the boundary square.}
The D.50 action is an automorphism of
\(\mathcal Z_h=\widehat B/R_{50}\).
Since it induces identity on \(G_3\), it preserves
the actual subgroup \(K_*/R_{50}\), the preimage
of \(P_*\) in this source. The inclusions
\(R_{50}\subset M_\infty\subset M_*\)
identify its maximal abelian pro-two kernel as
\(M_*/R_{50}\). This characteristic subgroup is
therefore preserved as well. Separately, the
retained map to \(\mathcal W_h\) makes
\(R_h/R_{50}\) invariant. Thus their intersection
\((R_h\cap M_*)/R_{50}\) is invariant, and the
action descends to
\(\mathcal Z_{h,*}=\widehat B/(R_h\cap M_*)\).
The quotient map
\(\mathcal Z_{h,*}\to\mathcal W_h\) remains equivariant.
The action also descends to \(\widehat B/M_*\).
These are the actual subgroup and quotient maps
defining \(\alpha_{0,*}\) and its inherited
conjugators in (D.56.10).

To identify the image of \(R_h\), consider the image
of \(C_h\) in \(\widehat B/M_*\).
It is normal, projects onto \(A_*\), and contains
\([A_*,U_*]=U_*^-\).
Modulo \(U_*^-\), its pro-two kernel commutes with
its odd \(A_*\)-quotient, and the extension splits
by the coprime averaging argument. Its image there
is pronilpotent.
The universal pronilpotent quotient of \(C_h\)
therefore kills exactly the moving part in this
image. This proves
\(\operatorname{im}(R_h)=U_*^-\).
One can verify the reverse inclusion directly at
every finite level: in a pronilpotent image the
three-primary image of \(C_h\) commutes with the
normal two-primary kernel, whereas their
commutators generate \(U_*^-\).

Apply the class-two construction of D.51 to these
actual subgroups. The odd averaging gives an
\(A_*\)-equivariant section in the affine space
of class-two lifts. Its full moving abelian image
is \(W_*\); commutators with the whole kernel give
exactly \(W_*\wedge V_*\).
The mixed invariant quotient is zero, and
fourth-power corrections remove its abelian
coordinate. This gives
\(\operatorname{im}(R_{50,*})=Z_*\).
It proves both the common quotient and the two
small projection kernels in (D.56.12).

The inclusions \(K_3\subset K_*\) imply
\(M_\infty\subset M_*\) and \(L_2\subset L_{2,*}\).
Taking the displayed intersections with \(R_h\)
gives the two full-source projection kernels
in (D.56.12). Thus these are quotient maps of
the same marked \(\widehat B\), with the literal
joint-source maps.

Choose the original class-two lift of the D.50
conjugators used to define \(q_{\alpha_0}\).
Its image defines the smaller lift.
Every pair correction and every central cusp
error maps by \(j_W\) and \(\theta\), respectively.
After taking the actual central correction
cokernels this gives the first equality in
(D.56.10). The same equality on local same-action
pairs gives
\(\tau\ell(\mathcal K)\subset\ell_*(\mathcal K_*)\).
It proves the entire reduced boundary square.
In particular
\[
 \ker\overline\tau
     ={\tau^{-1}(\ell_*(\mathcal K_*))
                              \over\ell(\mathcal K)}
 \tag{D.56.25}
\]
is its exact comparison kernel.
All of this retains the actual arithmetic seed
\(\alpha_0\); no value of that seed's remaining
invariant obstruction is inserted.

Finally, each \(G_*\)-element has \(|P_*|\) lifts
to \(G_3\). Equivariance of \(\theta\) therefore
proves (D.56.13). Since \(|P_*|\) is even,
that comparison is zero over \(\mathbb F_2\).
The boundary square, rather than a division by
this norm factor, supplies the evaluated
smaller arithmetic test.
\end{proof}
